\documentclass[12pt,a4paper]{article}
\usepackage[utf8]{inputenc}
\usepackage[russian]{babel}
\usepackage{geometry}
\geometry{a4paper, left=25mm, right=15mm, top=20mm, bottom=20mm}
\usepackage{hyperref}
\usepackage{listings}
\usepackage{color}

\definecolor{lightgray}{rgb}{0.95, 0.95, 0.95}
\lstset{
    backgroundcolor=\color{lightgray},
    basicstyle=\ttfamily\small,
    breaklines=true,
    frame=single,
    numbers=left,
    numbersep=5pt,
}

\title{Инструкция по подключению R-PRO (OPC UA Клиент) к удаленному OPC UA Серверу}
\author{Antigravity MES System}
\date{\today}

\begin{document}

\maketitle

\section{Введение}
В данном документе описывается процесс подключения среды имитационного моделирования R-PRO (выступающей в роли OPC UA Клиента) к нашему Java Spring Boot серверу (выступающему в роли OPC UA Сервера), когда они находятся на \textbf{разных физических компьютерах}.

Так как в конфигурации сервера (\texttt{OpcUaServerManager.java}) мы уже прописали прослушивание всех сетевых интерфейсов (\texttt{setBindAddress("0.0.0.0")}), сервер по умолчанию готов принимать внешние подключения.

\section{Шаг 1. Узнать IP-адрес компьютера с сервером}
На компьютере, где запущен Java OPC UA сервер, необходимо узнать его локальный IPv4 адрес.
\begin{enumerate}
    \item Откройте командную строку (\texttt{cmd}) или PowerShell.
    \item Введите команду:
    \begin{lstlisting}[language=bash]
    ipconfig
    \end{lstlisting}
    \item Найдите строку \texttt{IPv4 Address} (например, \texttt{192.168.1.50}). Запишите этот адрес.
\end{enumerate}

\section{Шаг 2. Настройка Брандмауэра Windows (Firewall)}
Чтобы сторонний компьютер мог подключиться к нашему серверу, необходимо открыть порт \texttt{4840} для входящих соединений на компьютере-сервере.
\begin{enumerate}
    \item Откройте \textbf{Брандмауэр Защитника Windows} (Windows Defender Firewall).
    \item Перейдите в \textbf{Дополнительные параметры} (Advanced settings).
    \item Выберите \textbf{Правила для входящих подключений} (Inbound Rules) и нажмите \textbf{Создать правило...} (New Rule).
    \item Выберите тип \textbf{Для порта} (Port) $\rightarrow$ Далее.
    \item Выберите \textbf{TCP} и укажите определенный локальный порт: \texttt{4840} $\rightarrow$ Далее.
    \item Выберите \textbf{Разрешить подключение} (Allow the connection) $\rightarrow$ Далее.
    \item Укажите профили (Доменный, Частный, Публичный) $\rightarrow$ Далее.
    \item Назовите правило, например, \texttt{OPC UA Server Port 4840}, и нажмите Готово.
\end{enumerate}

\section{Шаг 3. Подключение со стороны R-PRO}
Теперь необходимо настроить среду R-PRO на клиентском компьютере.
\begin{enumerate}
    \item Убедитесь, что оба компьютера находятся в одной локальной сети (подключены к одному Wi-Fi роутеру или соединены через VPN, например Radmin VPN / Hamachi).
    \item Откройте R-PRO на клиентском компьютере и перейдите в настройки \textbf{OPC UA} (раздел внешних подключений).
    \item В поле \textbf{Endpoint URL} (или Адрес сервера) введите:
    \begin{lstlisting}
    opc.tcp://[IP-адрес-сервера]:4840/
    \end{lstlisting}
    \textit{Пример: \texttt{opc.tcp://192.168.1.50:4840/}}
    \item Нажмите \textbf{Connect} (Подключиться). Индикатор должен загореться зеленым цветом.
    \item После успешного подключения вы сможете просматривать дерево тегов (Nodes) нашего сервера и привязывать переменные (например, \texttt{Robot1\_Command}) к 3D-моделям в сцене.
\end{enumerate}

\section{Решение возможных проблем}
\begin{itemize}
    \item \textbf{Связь прерывается или не устанавливается:} Попробуйте временно полностью отключить Брандмауэр на сервере. Если это помогло, значит порт \texttt{4840} был открыт неверно.
    \item \textbf{Компьютеры не видят друг друга:} Запустите команду \texttt{ping [IP-адрес-сервера]} на клиентском компьютере. Если ответа нет, проверьте сетевые настройки (например, видимость ПК в локальной сети Windows).
\end{itemize}

\end{document}
